\section{One checked owned-resource transfer}\label{sec:checked-owned-transfer}

Fix a source state, epoch and snapshot. A resource catalogue assigns each
identifier an owner, an asset and a positive integer amount. The state
records a finite marking \(M\) of active identifiers and a set \(U\) of
spent identifiers. A transfer signed by the catalogue owner of \(i\)
consumes \(i\) and creates a distinct identifier \(j\) of the same asset
and amount. The prechecks require \(M(i)=1\), \(M(j)=0\), and both
identifiers outside \(U\). The spend list is derived as \([i]\).

The verifier retains the original context and its cochain artifact. The
configured snapshot root must match the source state, and the actual
payload digest must match the artifact leaf. The supplied topology
authentication ties that snapshot to the state. These checks connect the
resource action to the configured cochain interpretation.

\paragraph{From the packet to the computed state.}
The two-slot manifest fixes the exact artifact and ciphertext before the
owner authorizes the transfer. The pre-proof draft records the source
parent and the roots of the source and computed successor state content.
The signed intent names that draft, the source evidence root, the original
public context, the transfer and the manifest root. Appending the signature
and payload proof gives the canonical packet. Its evidence digest and the
block digest complete the proposal. Both digest preimages exclude their own
digest fields.

The verifier parses the packet, confirms its canonical re-encoding, and
checks the exact expected metadata and intent. A different manifest root
therefore fails this check even on otherwise canonical bytes. The owner
signature is checked against the full intent and the payload proof against
the original interpretation. A profile requiring projection certificates
also checks the actual submitted bundle against its database, with every
required predecessor check on the same entries.

\begin{proposition}[Accepted update and replay rejection]
\label{prop:checked-owned-update}
For this fixed-context profile, suppose the packet checker accepts under
the configured authorization and exact-payload soundness contracts. The
accepted update uses the computed successor of the candidate constructed
from the parsed packet and original context. It preserves total value
for each asset, with
\[
  M'(i)=0,\qquad M'(j)=1,\qquad U'=U\cup\{i\}.
\]
If the initial inventory contains at most one copy per identifier and no
spent identifier is active, this invariant is preserved. The same action
cannot be accepted from this successor, even with a new proposal.
\end{proposition}

\begin{proof}
Acceptance supplies the resource and context prechecks and the exact parsed
packet checks. The signature contract authorizes its intent. The payload
contract gives exact acceptance under the original decoder using the checked
public-view equality. Together these facts derive candidate validity and
the accepted update. The single-copy flow removes \(i\) and inserts \(j\);
their matching catalogue entries give value preservation. All other
marking entries are retained, and the spend list adds \(i\). Freshness fails
for this same action in the resulting state, independently of the new
proposal.
\end{proof}

\paragraph{The recipient's view.}
The checked intent and computed history anchor retain the same manifest.
When its full ciphertext has the separate original-observation certificate,
the recipient checks the content ID on downloaded bytes and uses the
canonical issued token for its delivery scope. Under the hypotheses of
Proposition~\ref{prop:coupled-native-manifest}, this returns the original
coefficients or gives an accepted byte-hash collision. The catalogue value
in the state theorem comes from the resource record. An application in
which the encrypted coefficients determine the action supplies that
relation separately.

\paragraph{A finite reference execution.}
The runnable byte-owner instance starts with one copy of resource \(0\) and
no copy of resource \(1\). Both catalogue entries carry seven units of the
same asset, with different owners. It uses the existing exact zero payload,
canonical packet codec and fixture verifiers. The executable checker decides
resource predicates, compares canonical encodings of the public fields and
computes the finite successor. Lean proves that its acceptance and returned
state agree with the specified checker and successor on this profile.

The original canonical packet is \(3{,}524\) bytes. It moves the seven units
to resource \(1\) and marks resource \(0\) spent. The verifier rejects a
changed manifest root even with fixture authorization for the altered
intent. Both the original
packet and a freshly drafted attempt to spend \(0\) reject against the
returned state. The fixture uses a raw ciphertext slot and fixed verifier
data. Recipient recovery on genuine native bytes follows the separate
binding and delivery conditions above.
